% Source pins (SHA-256): PifoStatement.lean=fe9475ebfa2c462d27cb0a99781926074e1281396f2c9e8d32d6dfca788d814e; pifo-eta-decision.md=50159fe6aaf8273fc1d811f941de4e2246f7ec84e2270fefc1144c0a890e2ec9; pifo-elegant-v4.md=356ea38146ced437e20200db04e2ce220c73f9ac9fa892b746f783b8a2d49ab8; pifo-verify-all/VERIFICATION.md=c97ed5daa11c97f71a5d0e3867f96ddabe74438f1f598c049050c735a96b1887.
\section{An eta-style decision procedure}
\label{sec:eta-decision}

The decision procedure in \autoref{sec:deciding-equivalence} follows the
normalization route: it computes scheduler-independent pair and triple data
from which the canonical tree is determined.  There is also a direct,
pair-dependent route.  It recursively compares restrictions of the two input
schedulers and constructs only the meet preorders needed for that pair; it
never computes either scheduler's normal form.  The analogy with eta-equality
of functions is deliberately limited but useful: equality is established
extensionally, by comparing behavior at the arguments exposed by the two
objects, rather than by first normalizing their bodies.

Throughout this section \(A\) is a finite named alphabet, transported to
\(\Fin(k)\) when interfacing with the frozen statement, and all schedulers are
valid.  We retain the paper's notation \(F_S\) for the flush output with
\(\Some\) erased.  By \lemref{lem:balance}, every pop in a flush succeeds, so
inequality of these packet words is equivalent to inequality of the exact
\(\lean{List (Option (Fin k))}\) values returned by \(\run\).

Four operational facts support the recursion.  A run mentioning only
\(X\subseteq A\) is unchanged when performed in the restriction \(S|_X\), so
interleaved equivalence is hereditary under restriction.  A selector with one
active child is online-transparent.  Finally, an online-equivalent child may
be substituted below a fixed selector: the child receives projected pushes
and anonymous pops, and \lemref{lem:compression} accounts for gaps in its
global arrival stamps.  These facts preserve exact failed-pop behavior as
well as successful outputs.

At every recursive support \(X\) with \(|X|>1\), apply
\lemref{lem:decomposition} to expose
\begin{equation}
 \begin{aligned}
  S|_X&\intereq (P,r,(S_B)_{B\in P}),
  &S_B&\intereq S|_B,\\
  T|_X&\intereq (Q,s,(T_C)_{C\in Q}),
  &T_C&\intereq T|_C,
 \end{aligned}
 \label{eq:eta-exposures}
\end{equation}
where \(P,Q\) are proper partitions of \(X\).  Write
\(e_S(x,y)=\mathsf{eff}_S(x,y)\) for the effective comparison computed by
following the two assigned paths to their first decisive queue.  It can be
read directly from the finite paths or recovered from the two flushes \(xy\)
and \(yx\) by \autoref{tab:binary-recovery}; both arrival orders are needed to
distinguish a FIFO tie from either strict order.

\subsection{The pair-dependent local meet test}
\label{sec:eta-local}

Form the incidence meet
\begin{equation}
 R=P\wedge Q
   =\{B\cap C:B\in P,\ C\in Q,\ B\cap C\ne\varnothing\}.
 \label{eq:eta-meet}
\end{equation}
The local test first compares \(e_S(x,y)\) and \(e_T(x,y)\) for every pair in
distinct \(R\)-cells.  If they differ, one of \(xy\) and \(yx\) has different
flush output by \autoref{tab:binary-recovery}; the test runs both and returns the
one that differs.  Testing every cross-cell pair is essential.  On an
L-shaped triple one edge can be internal to a \(P\)-block and another internal
to a \(Q\)-block, so testing only pairs separated by both roots would assume
the very equivalence being decided.

Suppose all binary tests pass, and write \(e(x,y)\) for their common value.
Build a directed constraint graph using only pairs in distinct \(R\)-cells:
\begin{itemize}
\item a strict comparison \(x<_e y\) contributes a marked edge
      \(x\to y\);
\item a tie \(x\sim_e y\) contributes unmarked edges in both directions.
\end{itemize}
These constraints extend to a total preorder exactly when no strongly
connected component contains a marked edge.  In that case, collapse the
components, linearly extend the resulting DAG, and assign one natural rank to
each component.  Members of one component must receive the same rank: merely
placing them consecutively would incorrectly turn a required FIFO tie into a
strict comparison.  The resulting preorder \(t\) satisfies
\begin{equation}
 \begin{aligned}
  \cmp_t(x,y)&=\cmp_r(x,y)
    &&\text{if }P(x)\ne P(y),\\
  \cmp_t(x,y)&=\cmp_s(x,y)
    &&\text{if }Q(x)\ne Q(y).
 \end{aligned}
 \label{eq:eta-root-compatibility}
\end{equation}
Indeed, a \(P\)-separated pair is decided at \(S\)'s exposed root and a
\(Q\)-separated pair at \(T\)'s; the preceding binary phase makes the two
prescriptions agree whenever both roots separate the pair.

If a marked edge lies in an SCC, the cycle itself is only a graph failure, not
yet an observational certificate.  The shortest-marked-cycle argument behind
\lemref{lem:meet-preorder} has the following constructive adaptation.  Shorten
a marked closed walk whenever vertices two steps apart occupy distinct cells
and their three-cell comparisons are transitive.  A failed shortcut exposes
three distinct cells on which \(e\) is not a preorder.  If no shortcut and no
such triple existed, the walk would alternate between two cells.  All
comparisons between those fixed cells come from one genuine root preorder---
from \(r\) if their \(P\)-blocks differ, and otherwise from \(s\)---so a
marked alternating cycle is impossible.

The resulting bad triple must form an L in the \(P\times Q\) incidence grid.
If all three \(P\)-coordinates were distinct, \(r\) would govern every edge;
if all three \(Q\)-coordinates were distinct, \(s\) would.  Nontransitivity
therefore gives distinct \(a,b,c\) with
\begin{equation}
  a\le_e b,\qquad b\le_e c,\qquad c<_e a.
 \label{eq:eta-bad-triple}
\end{equation}
Use the arrival word \(\sigma=abc\), whose order orients either of the first
two ties forward.  Let \(\prec_\sigma\) denote the resulting stable pairwise
precedence.  After possibly exchanging \((S,P)\) with \((T,Q)\), cyclically
name the L's common corner \(x\) so that
\(x\prec_\sigma y\prec_\sigma z\prec_\sigma x\), without changing
\(\sigma\).  The possible exchange ensures that \(y\), rather than \(z\), is
the neighbor sharing \(x\)'s \(P\)-block.  Thus the \(P\)-decomposition groups
\(\{x,y\}\mid\{z\}\): its root slots are
\(\{x,y\},z,\{x,y\}\), its pair child emits \(x,y\), and its output is
\(xzy\).  The \(Q\)-decomposition groups
\(\{x,z\}\mid\{y\}\): its slots are
\(\{x,z\},y,\{x,z\}\), its pair child emits \(z,x\), and its output is
\(zyx\).  This is the L-diamond calculation of \lemref{lem:diamond}, now used
after actual binary agreement rather than under an assumed common flush
transformer.  An implementation need not reconstruct the orientation: it can
run all six permutations of the extracted triple and return the first unequal
one.

\begin{lemma}[Local extension or witness]
\label{lem:eta-local-alternative}
At a recursive support, the finite local test returns exactly one of:
\begin{enumerate}
\item a total preorder \(t\) satisfying
      \autoref{eq:eta-root-compatibility}; or
\item an injective packet word \(w\), with \(|w|=2\) or \(3\), such that
      \(F_S(w)\ne F_T(w)\).
\end{enumerate}
\end{lemma}

\begin{proof}
A binary mismatch gives the second alternative on one of the two arrival
orders.  After binary agreement, the SCC criterion either constructs \(t\) or
the marked-walk shortening produces the L-diamond witness above.  These are
the exhaustive, mutually exclusive branches of the local test.
\end{proof}

\subsection{Recursive comparison without normal forms}

The procedure \(\mathsf{Compare}(X)\) always compares the two original input
schedulers restricted to \(X\); equal subsets are memoized.
\begin{enumerate}
\item If \(|X|\le1\), return an equivalence base certificate.
\item Otherwise expose \(S|_X\) and \(T|_X\) as in
      \autoref{eq:eta-exposures}, form \(R=P\wedge Q\), and run the local test.
\item Return immediately if the local test found a distinguishing word.
      Otherwise retain its preorder \(t\).
\item Recursively compare every distinct block occurring in \(P\cup Q\).
      A word returned by a child remains a witness at \(X\) by restriction
      exactness.
\item If all children succeed, return a proof node recording \(P,Q,R,t\) and
      references to their certificates.
\end{enumerate}
Both top partitions are proper, so every recursive block is smaller than
\(X\).  Strong recursion therefore terminates even without memoization; the
memo table is needed for the polynomial implementation analyzed below.

Equivalently, for \(|X|>1\) the procedure establishes the exact recursive
characterization
\begin{equation}
\boxed{
 S|_X\intereq T|_X
 \quad\Longleftrightarrow\quad
 \begin{array}{l}
 e_S=e_T\text{ on pairs in distinct }P\wedge Q\text{-cells},\\
 \text{those comparisons extend to a total preorder }t,\\
 S|_Y\intereq T|_Y\text{ for every distinct }Y\in P\cup Q.
 \end{array}}
\label{eq:eta-recursive-characterization}
\end{equation}
The partition, extension, and proof bridge are chosen for this pair at this
recursive support.  None is a normal form for either input.

\subsection{The refinement bridge}
\label{sec:eta-bridge}

The successful local node is sufficient because it supports a one-sided
refinement that does not assume the two schedulers equivalent.  For every
\(D\in R\), choose
\begin{equation}
 U_D=S|_D,
 \qquad
 U=(R,t,(U_D)_{D\in R}).
 \label{eq:eta-bridge-U}
\end{equation}
For \(C\in Q\), let \(W_C=U|_C\): after inactive children are removed, this
is a \(t|_C\)-root over the nonempty cells \(D=B\cap C\), with child \(U_D\).

\begin{lemma}[One-sided refinement]
\label{lem:eta-one-sided-refinement}
If the local test returns \(t\), then, for every \(C\in Q\),
\begin{equation}
  S|_C\intereq W_C.
 \label{eq:eta-one-sided-refinement}
\end{equation}
This conclusion does not assume \(S\intereq T\).
\end{lemma}

\begin{proof}
Restrict the exposed \(P\)-root of \(S\) to \(C\).  Its active colors are
exactly the nonempty cells \(D=B\cap C\).  Below color \(D\), its child is
\(S_B|_D\).  The exposure equivalence in
\autoref{eq:eta-exposures}, heredity, and restriction idempotence give
\begin{equation}
  S_B|_D\intereq (S|_B)|_D=S|_D=U_D.
 \label{eq:eta-cell-child}
\end{equation}
This also covers the unary-prefix and terminal-leaf-expansion cases used to
obtain the exposure; unary transparency alone would not justify the displayed
child equivalence.

Two distinct cells inside the fixed \(Q\)-block \(C\) lie in distinct
\(P\)-blocks.  Hence \autoref{eq:eta-root-compatibility} says that \(r\) and
\(t\) have the same exact comparison between their colors, including ties.
Apply \lemref{lem:colored} and then substitute the equivalent children from
\autoref{eq:eta-cell-child}.  The resulting scheduler is \(W_C\), proving
\autoref{eq:eta-one-sided-refinement}.  Balance gives common \(\None\) outputs
on over-pops.
\end{proof}

\begin{lemma}[Successful-node sufficiency]
\label{lem:eta-successful-node}
Suppose the local test returns \(t\) and recursive certificates prove
\begin{equation}
 E_Y:S|_Y\intereq T|_Y
 \qquad\text{for every distinct }Y\in P\cup Q.
 \label{eq:eta-child-certificates}
\end{equation}
Then \(S|_X\intereq T|_X\).
\end{lemma}

\begin{proof}
First transform \(T|_X\) into the bridge \(U\).  For each \(C\in Q\), the
exposure, recursive certificate, and one-sided refinement give
\begin{equation}
 T_C\intereq T|_C\intereq S|_C\intereq W_C.
 \label{eq:eta-T-child-chain}
\end{equation}
Substitute \(W_C\) below the exposed \(s\)-root.  This leaves an outer
\(s\)-selector over \(Q\) and, inside each block, a \(t\)-selector over the
\(R\)-cells.  By \autoref{eq:eta-root-compatibility} and
\lemref{lem:colored}, recolor the outer selector from \(s\) to \(t\) without
changing its \(Q\)-trace.  The two selector levels now use identical
\((t,\arr)\) keys and collapse occurrence-for-occurrence as in
\autoref{sec:induction}.  Thus
\begin{equation}
  T|_X\intereq U.
 \label{eq:eta-T-to-U}
\end{equation}

For each \(B\in P\), restrict \autoref{eq:eta-T-to-U} to \(B\) and use the
recursive certificate:
\begin{equation}
 S_B\intereq S|_B\intereq T|_B\intereq U|_B.
 \label{eq:eta-S-child-chain}
\end{equation}
After inactive children are deleted, \(U|_B\) is the \(t|_B\)-root over the
cells \(D\subseteq B\).  Substitute these restrictions below \(S\)'s exposed
\(r\)-root, recolor \(r\) to \(t\), and collapse the duplicated \(t\)-levels.
This gives \(S|_X\intereq U\), which combines with
\autoref{eq:eta-T-to-U} to prove the claim.

The collapse couples pending source occurrences, not merely packet labels:
the outer queue, the relevant inner queue, and the flat queue choose the same
unique \((t,\arr)\)-minimum.  It therefore covers FIFO ties and repeated
packet types.  Balance makes selected children nonempty; a common over-pop
returns \(\None\) and changes neither state.
\end{proof}

Necessity in \autoref{eq:eta-recursive-characterization} is direct.  If the two
restrictions are interleaved-equivalent, heredity supplies every block
equivalence, their two-type restrictions agree, and a marked SCC is
impossible because \lemref{lem:eta-local-alternative} would replay an unequal
three-push flush.  This argument does not reconstruct either top partition
from observations.

The successful-node proof also does not derive refined child transformers from
a globally common flush function.  Recursive certificates directly supply
equivalence on every original \(P\)- and \(Q\)-block, while
\lemref{lem:eta-one-sided-refinement} supplies the missing bridge.  Thus a full
cell-autonomy or meet control-word calculation for constructing the refined
children is not an implicit premise of this comparator.

\subsection{Correctness and small counterexamples}

\begin{theorem}[Eta-style comparison]
\label{thm:eta-comparison}
For every finite \(X\), \(\mathsf{Compare}(X)\) terminates and returns exactly
one of:
\begin{itemize}
\item \(\mathsf{Equivalent}(c)\), where \(c\) proves
      \(S|_X\intereq T|_X\); or
\item \(\mathsf{Distinguishing}(w)\), where \(w\) is injective,
      \(|w|\in\{2,3\}\), and \(F_S(w)\ne F_T(w)\).
\end{itemize}
In particular,
\begin{equation}
 \mathsf{Compare}(X)=\mathsf{Equivalent}
 \quad\Longleftrightarrow\quad
 S|_X\intereq T|_X.
 \label{eq:eta-comparison-correct}
\end{equation}
\end{theorem}

\begin{proof}
Use joint strong induction on \(|X|\), proving termination and the stated
semantic alternative simultaneously.  For \(|X|=0\), every operation is a
failed pop.  For \(|X|=1\), every successful pop emits the unique type and
success depends only on the pending count.

For \(|X|>1\), a binary mismatch is a correct two-letter witness by
\autoref{tab:binary-recovery}; a graph failure is a correct three-letter witness
by the L-diamond branch of \lemref{lem:eta-local-alternative}.  A word returned
from a proper block remains unequal in the original restrictions because it
contains no outside pushes, so even its arrival stamps are unchanged.  If all
recursive calls succeed, \lemref{lem:eta-successful-node} constructs the
equivalence proof.  Every call is on a smaller support, completing the joint
induction.  Thus no equivalent pair can take a negative branch, since every
such branch contains a replayed unequal run, and every positive branch has the
bridge proof above.
\end{proof}

Let
\begin{equation}
 \mathcal W_3(A)=
 \{w\in A^*:w\text{ is injective and }|w|\in\{2,3\}\}.
 \label{eq:eta-W3}
\end{equation}

\begin{corollary}[Two- or three-push counterexamples]
\label{cor:eta-small-counterexample}
For valid stateless schedulers over \(A\),
\begin{equation}
 S\intereq T
 \quad\Longleftrightarrow\quad
 \forall w\in\mathcal W_3(A).\ F_S(w)=F_T(w).
 \label{eq:eta-finite-observation}
\end{equation}
Consequently, if \(S\not\intereq T\), there is a differing flush word with
exactly two or three pushes of pairwise-distinct packet types.
\end{corollary}

\begin{proof}
The forward implication is immediate.  For the reverse implication, apply
the induction in \autoref{thm:eta-comparison}.  Binary equality supplies the
common cross-meet comparisons; if they did not extend, the L-diamond would
give a word in \(\mathcal W_3(A)\).  Every proper \(P\)- or \(Q\)-block
inherits the same hypothesis, so its recursive comparison succeeds, and
\lemref{lem:eta-successful-node} finishes.  Equivalently, run the comparator:
its only negative leaves are words in \(\mathcal W_3(A)\).
\end{proof}

A simple streaming implementation may therefore test
\begin{equation}
 2\binom{k}{2}+6\binom{k}{3}=k(k-1)^2
 \label{eq:eta-stream-count}
\end{equation}
words and stop on the first difference.  This observational implementation
does not materialize a profile, but its correctness proof is the recursive
meet argument above.  The structural comparator can do less work on shallow
inputs.

\subsection{Memoized complexity}
\label{sec:eta-complexity}

A machine bound requires a representation for the mathematical assignment
function.  Assume the topology and all assigned paths are finite encoded data;
let \(N\) be their read and tabulation cost, measured in words for a word-RAM
bound and in bits for a bit-cost bound.  Validity is promised or checked during
preprocessing.  Type and cluster identifiers and compressed rank classes fit
one word, and \(b\) is the maximum bit length of an input rank when bit cost is
charged.  Thus inactive topology and long common unary paths are included in
\(N\), rather than treated as free.  The asymptotic discussion below assumes
\(k\ge2\); the cases \(k=0,1\) terminate at the base case after the input is
read.

For each input scheduler, contract active unary nodes and apply the terminal
leaf expansion from \lemref{lem:decomposition} recursively.  The result is an
\emph{active support hierarchy}: a rooted laminar tree with the \(k\) types as
singleton leaves and no unary internal node.  It has at most \(k-1\) internal
nodes for \(k\ge1\).  This is syntax-directed indexing of the input, not a
semantic canonical form.  At every active selector, retain the rank class of
each descendant type.  For input \(i\in\{1,2\}\), put
\begin{equation}
 a_i(x)=\#\{\text{active support clusters on \(x\)'s ancestor chain}\},
 \qquad
 L_i=\sum_x a_i(x).
 \label{eq:eta-L}
\end{equation}
Up to the harmless convention for singleton clusters,
\begin{equation}
  k\le L_i\le k^2.
 \label{eq:eta-L-range}
\end{equation}
Rank compression at the selectors costs
\(O(bL_i\log k)\) bit operations by comparison sorting.  LCA and
selector/type lookup tables then make an effective two-type comparison
constant-time.

Every subset reached by the recursion has the form
\begin{equation}
  X'=X_1\cap X_2,
 \label{eq:eta-cluster-intersection}
\end{equation}
where \(X_i\) is a cluster in input \(i\)'s active hierarchy.  Indeed, the
exposed top of \(S|_{X'}\) is the least first-hierarchy cluster containing
\(X'\), and its blocks are intersections of \(X'\) with that cluster's
children; the second hierarchy is symmetric.  Starting from the full
alphabet, recursion on a \(P\)- or \(Q\)-block therefore preserves the
intersection form.

Replacing \(X_1,X_2\) by the least clusters containing \(X'\) gives a
canonical memo key and preserves the intersection.  If
\(X'=X_1^0\cap X_2^0\), the two least clusters satisfy
\(X_i\subseteq X_i^0\) while still containing \(X'\), so their intersection
is exactly \(X'\).  Each hierarchy has \(O(k)\) clusters, so there are only
\(O(k^2)\) nonempty memo states.  Memoization is crucial:
opposite combs can expose two different \((m-1)\)-blocks and make the unfolded
recurrence satisfy
\begin{equation}
  T(m)\ge2T(m-1)+O(m),
 \label{eq:eta-unmemoized}
\end{equation}
so the literal recursive call tree can be exponential even for equivalent
inputs.

Scanning every pair independently in every state gives the safe but loose
bound \(O(N+k^4)\).  Instead enumerate exactly the pairs separated by the
current \(P\) or \(Q\), which are the pairs in distinct meet cells.  In the
sums below, \(x<y\) uses the fixed alphabet enumeration only to select each
unordered pair once.  Fix types \(x,y\).  At a state where \(P\) separates
them, the first hierarchy's
cluster coordinate is forced to \(\operatorname{lca}_1(x,y)\); the other
coordinate can be only a common ancestor of \(x,y\), giving at most
\(\min(a_2(x),a_2(y))\) choices.  The symmetric charge applies when \(Q\)
separates them.  Hence the total number of generated pair constraints is
\begin{equation}
\begin{aligned}
 E_{\mathrm{all}}
 &=O\!\left(
   \sum_{x<y}\min(a_1(x),a_1(y))
   +\sum_{x<y}\min(a_2(x),a_2(y))\right)\\
 &=O\bigl(k(L_1+L_2)\bigr).
\end{aligned}
\label{eq:eta-edge-charge}
\end{equation}
The last step uses \(\min(u,v)\le(u+v)/2\); in the unordered-pair sum each
\(a_i(x)\) occurs \(k-1\) times.  The total number of type occurrences over
all memo states also satisfies
\begin{equation}
 Z\le\sum_x a_1(x)a_2(x)
   \le k\min(L_1,L_2).
 \label{eq:eta-vertex-charge}
\end{equation}
Partition construction, pair checking, SCCs, and topological extension are
linear in these state vertices and edges.  Shortening a marked cycle to its
bad triple costs \(O(|X'|^2)\) within the single failing state and remains
within the same aggregate bound.  A simpler implementation may enumerate
triples instead; this adds \(O(k^3)\), preserves the cubic worst-case bound,
but can lose the sharper quadratic bound on a shallow failing input.

\begin{theorem}[Eta comparator complexity]
\label{thm:eta-complexity}
After input and rank preprocessing, the optimized memoized comparator uses
\begin{equation}
 O\bigl(k(L_1+L_2)\bigr)\subseteq O(k^3)
 \label{eq:eta-word-time}
\end{equation}
word-RAM operations.  Its working decision space, including the encoded
input but excluding a retained positive certificate, can be \(O(N+k^2)\) by
storing the memo keys and one local graph at a time.
\end{theorem}

The word bound assumes that an identifier, graph endpoint, memo key, and
compressed rank class occupy \(O(\log k)\) bits.  Charging their bit
operations and the comparison sorting gives the conservative combined bound
\begin{equation}
 O\bigl(
   N+b(L_1+L_2)\log k
    +k(L_1+L_2)\log k
 \bigr),
 \label{eq:eta-bit-time}
\end{equation}
where \(N\) is measured in bits.  In particular, the structural word term in
\autoref{eq:eta-word-time} acquires a \(\log k\) factor; unbounded natural ranks
are not silently unit-cost.  When both hierarchies are shallow,
\(L_1+L_2=\Theta(k)\), the direct structural work is \(O(k^2)\).  Opposite
combs can make \(L_i=\Theta(k^2)\) and attain the cubic charge.

This adaptive improvement is not a subcubic worst-case claim.  For comparison,
one can store one three-valued code for every unordered pair and six flush
outputs, one for each input permutation of every unordered triple, for
\begin{equation}
 \binom{k}{2}+6\binom{k}{3}=\Theta(k^3)
 \label{eq:eta-dense-profile-size}
\end{equation}
constant-size codes.  After the same hierarchy and LCA preprocessing, a
two- or three-type restriction has only constantly many active branching
selectors, reached by jumping to the LCA of the selected types and, when
needed, of the remaining pair.  Its drain therefore uses \(O(1)\)
compressed-rank comparisons.  The scheduler-independent profile takes
\(O(k^3)\) additional time to write and is reusable in later comparisons.
The direct comparator has the same cubic
worst-case exponent, but avoids materializing two dense profiles, may stop at
the first witness, and can take only quadratic structural work on shallow
pairs.  The \(\Omega(k^3)\) output cost applies only when that explicit dense
profile is written, and the opposite-comb charge analyzes this implementation;
neither is a lower bound for every exact algorithm or canonical syntax.  In
particular, no claim is made that worst-case \(o(k^3)\) decision or
representation after input reading and rank compression is impossible.

\subsection{Certificates in both directions}

On inequivalence, the returned \(w\) is injective and has two or three
letters.  The certificate is the operation word \(\flushOps(w)\), of length
four or six.  A checker simply runs both original schedulers and verifies that
their exact \(\lean{List (Option (Fin k))}\) outputs differ.  Once replayed,
this certificate is independent of the comparator's correctness proof.  Its
constant operation length follows because primitive failures are binary-table
failures or L-diamonds and restriction lifting never changes the word.  Its
two or three type identifiers still occupy \(O(\log k)\) bits each.

There is no single positive test word.  An equivalence certificate is a
finite, pair-dependent proof DAG keyed by the reachable cluster intersections
in \autoref{eq:eta-cluster-intersection}.  Every non-base node records
\begin{itemize}
\item its partitions \(P,Q,R\);
\item rank classes for \(t\), preserving every required cross-cell strict
      comparison and tie; and
\item references to certificate nodes for every distinct block in
      \(P\cup Q\).
\end{itemize}
The checker recomputes both exposures, checks the recorded
\(R=P\wedge Q\), and recomputes both effective comparisons for every required
cross-cell pair.  It checks equality, verifies \(t\) against the common
strict-or-tie value, checks that every required proper block is referenced,
and applies the fixed
one-sided-refinement, substitution, recoloring, and collapse inference from
\autoref{sec:eta-bridge}.  Supplying \(t\) avoids constructing an extension but
not checking its constraints.  It scans the same \(Z\) occurrences and
\(E_{\mathrm{all}}\) pair constraints as successful decision, so its checking
cost after preprocessing is
\begin{equation}
 O(Z+E_{\mathrm{all}})=O\bigl(k(L_1+L_2)\bigr)
 \label{eq:eta-certificate-check}
\end{equation}
word operations, cubic for opposite combs.  Storing every chosen \(t\) and
the DAG metadata takes \(O(Z\log k)\) bits, at most \(O(k^3\log k)\).

The two certificate directions have different modes of checking.  The negative
word is independently replayable against the frozen \(\run\) definition.
Development~16 of the pinned verification manifest encodes the comparator as a
terminating computable definition and kernel-checks its soundness,
completeness, and two-or-three-push witness-size theorem; its soundness and
witness-size gates use only \(\{\lean{propext},\lean{Quot.sound}\}\).  Its
trust-level-zero audit passes, and its \lean{\#eval} probes execute the
comparator and replay separating witnesses through the frozen semantics.  In
the positive direction, \lean{etaCompare\_complete} validates the comparator's
equivalent outcome.  A positive proof DAG cannot be validated merely by
replaying scheduler outputs; its recursive inference is instead covered by
that kernel-checked completeness theorem.  A common dense profile of
\autoref{eq:eta-dense-profile-size}, together with replay of all its entries and
\corref{cor:eta-small-counterexample}, is an alternative canonical positive
certificate.  The proof DAG, like the comparator itself, remains pair-dependent.
This is the precise eta-style advantage: it decides
equivalence without first computing normal forms, while making no claim of a
better worst-case exponent.
